% Target: ~200-250 words, single paragraph. Rewrite last, once results exist.
The design of the superconducting magnet system of a fusion device is usually
carried out downstream of the plasma and magnetic configuration studies, and
iterated by hand between conductor, winding-pack and structural analyses. We
present \MADE{} (Magnet Design Explorer), a code that closes this gap by taking
the \emph{target magnetic configuration} as input --- field on axis, coil
number and ripple for the toroidal field, flux swing for the central solenoid,
coil currents and positions from the equilibrium for the poloidal field, the
axial field profile for a magnetic mirror --- and returning the set of
feasible magnet designs that realise it. \MADE{} is organised in coil-agnostic
modules: (i) conductor sizing (superconductor selection among NbTi, Nb$_3$Sn
and REBCO, cable cross-section, copper for quench protection via hot-spot
analysis), (ii) winding-pack grading, and (iii) structural sizing against
Tresca allowables. These modules are combined into coil-specific solvers for
toroidal field (TF), central solenoid (CS) and poloidal field (PF) coils, and
--- since a mirror machine is a series of solenoids and rings --- the same
modules design mirror magnets without new physics. Rather than a single
optimum, \MADE{} exhaustively explores the design space and returns every
feasible design point, allowing trade-off studies. \todo{verification result,
e.g. agreement with reference designs within X\%}. Finally, we lay out the
extension to stellarators, where non-planar coils lose the wedged support of
tokamak TF coils: the conductor is sized as in the TF solver without the
wedge, and the structural verification is delegated to a global beam--shell
model able to follow non-planar centrelines, which also generalises the
analysis of D-shaped TF coils.
